The structure equations below are those of the derived metric of the gravitational field, whose unique torsion-free solution is its rotation-gauge field written in the coframe (Fact \claimref{synthesis.imported-gtg-dictionary}(ii)); the matrix $\Omega$ is the coframe's logarithmic rate of the statement. (i) The spatial coframe is $e = \Delta G\theta_{0}$ as a column of one-forms, so $\partial_{\tau}e = (\partial_{\tau}\Delta)\Delta^{-1}e + \Delta(\partial_{\tau}G)G^{-1}\Delta^{-1}e = \tau^{-1}(P + \Delta\Omega\Delta^{-1})e = \tau^{-1}Be$, since $\partial_{\tau}\tau^{p_{l}} = (p_{l}/\tau)\tau^{p_{l}}$. Hence $de^{l} = \tau^{-1}B_{lm}\,e^{0}\wedge e^{m} + d_{x}e^{l}$, with $d_{x}$ the exterior derivative along the slice at fixed $\tau$, which has no $e^{0}$ leg; and $de^{0} = 0$. In Cartan's first structure equation the $e^{0}\wedge e^{m}$ part of $de^{l}$ is matched by the boost block $\omega^{l}{}_{0} = K_{lm}e^{m}$, which is symmetric in $(l, m)$, and by the $e^{0}$-leg of the rotation block, which is antisymmetric; so the boost block is $\tau^{-1}\mathrm{sym}(B)$ and the $e^{0}$-leg of the rotation block is $\tau^{-1}$ times the antisymmetric part of $B$, the sign fixed by the convention $\omega^{l}{}_{m} = -\omega^{lm}$ of the diagonal Lemma \claimref{dark-energy.variable-exponents-residue-pointwise-corrections-are}, whose boost block $(p_{l}/\tau)e^{l}$ is the case $G = 1$, $B = P$. Independently, with $h = e^{\mathsf T}e$ the spatial metric in coordinates, $K = \tfrac12 h^{-1}\partial_{\tau}h = \tfrac12 e^{-1}\bigl((\partial_{\tau}e\,e^{-1})^{\mathsf T} + \partial_{\tau}e\,e^{-1}\bigr)e$, so $\tau K$ is similar to $\mathrm{sym}(B)$. The script checks the last statement exactly from the coordinate metric $h = \theta_{0}^{\mathsf T}G^{\mathsf T}\Delta^{2}G\theta_{0}$ at the Kasner point $p = (-\tfrac27, \tfrac37, \tfrac67)$, with $\tau = t^{7}$ so that every power is rational, exact derivatives by dual numbers, a non-diagonal constant $\theta_{0}$, and $G = \exp(c\log\tau\,N)$ for $N$ triangular in either direction, $\log\tau$ entering as an independent symbol at two rational values (the identity is polynomial in it): twelve cases, characteristic polynomials equal; and for $G = 1$ the characteristic polynomial of $\tau K$ is that of $P$, $(1, 0, -\tfrac{36}{343})$.

(ii) $P$ is constant, so $\tau^{-1}\mathrm{sym}(B)$ and $\tau^{-1}\mathrm{antisym}(B)$ are $O(\tau^{-1})$ if and only if $B$, equivalently $\Delta\Omega\Delta^{-1}$, is bounded, and their residues are $P$ and $0$ if and only if $\Delta\Omega\Delta^{-1} \to 0$, whose entries are $\Omega_{lm}\tau^{p_{l} - p_{m}}$. If $\Omega$ is bounded and $\Omega_{lm} \neq 0$ only where $p_{l} > p_{m}$, every nonzero off-diagonal entry carries a positive power of $\tau$ and tends to zero, and the diagonal of a strictly triangular $\Omega$ vanishes, so the residue is unchanged. For $G = \exp(c\log\tau\,N)$, $N$ commutes with $G$, so $\Omega = \tau\cdot c\tau^{-1}N = cN$; for $N$ strictly triangular in some order each unordered pair $\{l, m\}$ carries at most one of $N_{lm}$, $N_{ml}$, the diagonal of $\mathrm{sym}(B)$ is $P$, its $(l,m)$ and $(m,l)$ entries are $\tfrac{c}{2}N_{lm}\tau^{p_{l} - p_{m}}$, and $\mathrm{tr}((\tau K)^{2}) = \mathrm{tr}(\mathrm{sym}(B)^{2})$ is the displayed sum. A term with $p_{l} < p_{m}$ grows without bound, and $\mathrm{tr}((\tau K)^{2}) = \tau^{2}K^{i}{}_{j}K^{j}{}_{i}$ is a scalar, so no change of frame removes the growth. The script pins the sums at $t = \tfrac12, \tfrac14, \tfrac18, \tfrac1{16}$ for one $N$ of each kind ($\tfrac{821657}{819200}$ decreasing to $1$, and $\tfrac{1184569}{25}$ increasing).

(iii) For $G = 1 + fE_{12}$ with $f = \tau^{3/7}$, $G^{-1} = 1 - fE_{12}$ and $\Omega = \tau f'E_{12} = \tfrac37\tau^{3/7}E_{12}$, which tends to zero (the frame converges to the identity at rate $f' = O(\tau^{-4/7}) = o(1/\tau)$), while $(\Delta\Omega\Delta^{-1})_{12} = \tfrac37\tau^{3/7 + p_{1} - p_{2}} = \tfrac37\tau^{-2/7}$ is unbounded, $p_{1} < p_{2}$. By (ii), $\mathrm{tr}((\tau K)^{2}) = 1 + 2\bigl(\tfrac{3}{14}\tau^{-2/7}\bigr)^{2} = 1 + \tfrac{9}{98}\tau^{-4/7}$; the script computes it from the coordinate metric at four values of $t$ ($\tfrac{121}{49}, \tfrac{1201}{49}, \tfrac{18481}{49}, \tfrac{294961}{49}$).

(iv) In the synchronous gauge the $00$ component of the Ricci tensor is $-\partial_{\tau}\mathrm{tr}K - \mathrm{tr}(K^{2})$ and the Hamiltonian constraint is $(\mathrm{tr}K)^{2} - \mathrm{tr}(K^{2}) + R^{(3)} = 2\Lambda$ up to the sign convention for $R^{(3)}$ (standard; cited from memory). With $\tau K = R + o(1)$ and $\tau\partial_{\tau}(\tau K) \to 0$, $\partial_{\tau}\mathrm{tr}K = -\tau^{-2}\mathrm{tr}R + o(\tau^{-2})$, so $R_{00} = \Lambda$ gives $\mathrm{tr}R - \mathrm{tr}R^{2} = O(\tau^{2})$, hence $\mathrm{tr}R = \mathrm{tr}R^{2}$; with $\tau^{2}R^{(3)} \to 0$ the constraint gives $(\mathrm{tr}R)^{2} = \mathrm{tr}R^{2}$. So $\mathrm{tr}R \in \{0, 1\}$, and for $\mathrm{tr}R > 0$ the real eigenvalues of the symmetric $R$ satisfy $\sum = \sum(\cdot)^{2} = 1$, the Kasner circle; $\Lambda$ enters only at $O(\tau^{2})$, as in Proposition \claimref{dark-energy.cosmological-constant-does-not-change-residue}. The relations are the leading field equations of Proposition \claimref{dark-energy.frozen-point-datum-kasner-wall}(ii) read for a non-diagonal residue, and under the criterion of (ii) $R = P$. The script checks the residue algebra at the chosen point and a control off the circle. The criterion of (ii) is the symmetry-wall structure of the Iwasawa-frame analysis of cosmological billiards (Damour, Henneaux and Nicolai 2003; Heinzle, Uggla and Rohr 2009; cited from memory); the computation here is self-contained. The hypothesis $\tau^{2}R^{(3)} \to 0$ is not derived here, and it is not generic: on the Kasner circle some exponent is $\leq 0$, and for generic vacuum data with a negative exponent the slices' curvature is expected to dominate and produce a bounce (argued from the same literature; filed alongside, synthesis.unit-boost-residue-needs-nonpositive-exponent); in the diagonal class the momentum constraints of \claimref{synthesis.diagonal-kasner-momentum-constraint} are the next order, and their non-diagonal analogue is not treated.
